%%%%%%%%%%%%%%%%%%%%%%%%%%%%%%%%%%%%%%%%%%%%%%%%%%%%%%%%%%%%%%%%%%%%%%
%%                     Catalysis
%%%%%%%%%%%%%%%%%%%%%%%%%%%%%%%%%%%%%%%%%%%%%%%%%%%%%%%%%%%%%%%%%%%%%%
%\color{blue}
\subsection{Glyph: \glyph{Catalysis}}\label{sec:catalysis}

A catalysis is a particular case of stimulation, where the effector affects
positively the flux of a process represented by the target process. The positive effect on the process is due to the lowering of the activation energy of a reaction.

\begin{glyphDescription}
 \glyphSboTerm SBO:0000172 ! catalysis.
 \glyphOrigin Any \glyph{EPN} (\sect{EPNs}) or any \glyph{logical operator} (\sect{logic}).
 \glyphTarget Any \glyph{process node} (\sect{PNs}).
 \glyphNode The target extremity of a \glyph{catalysis} carries an empty circle.
 \end{glyphDescription}

\begin{figure}[H]
  \centering
  \includegraphics[scale = 0.5]{images/catalysis}
  \caption{The \PD glyph for \glyph{catalysis}.}
  \label{fig:catalysis}
\end{figure}


